\documentclass[10pt,a4paper]{article}

% ---------- fonts & language ----------
\usepackage{fontspec}
\setmainfont{DejaVu Sans}
\usepackage{polyglossia}
\setmainlanguage{german}
\usepackage{microtype}

% ---------- layout ----------
\usepackage[a4paper,top=1.6cm,bottom=1.6cm,left=1.9cm,right=1.9cm]{geometry}
\setlength{\parindent}{0pt}
\setlength{\parskip}{1pt}
\pagestyle{empty}

% ---------- color & boxes ----------
\usepackage[table]{xcolor}
\definecolor{navy}{RGB}{23,41,79}
\definecolor{navyfaint}{RGB}{224,229,238}
\definecolor{gray1}{RGB}{95,95,95}
\definecolor{gray2}{RGB}{243,243,244}
\usepackage[skins]{tcolorbox}

% ---------- tables & lists ----------
\usepackage{array}
\usepackage{tabularx}
\usepackage{enumitem}
\setlist[itemize]{topsep=0.8pt,itemsep=0.8pt,parsep=0pt,leftmargin=1.15em,label=\textcolor{navy}{\textbullet}}

% ---------- section headings ----------
\usepackage{titlesec}
\titleformat{\section}
  {\normalfont\Large\bfseries\color{navy}}{}{0pt}{}
  [{\color{navy}\titlerule[1.1pt]}]
\titlespacing*{\section}{0pt}{8pt}{5pt}

% ---------- links ----------
\usepackage{hyperref}
\hypersetup{
  colorlinks=true, linkcolor=navy, urlcolor=navy, citecolor=navy,
  pdftitle={Lebenslauf Iurii Pigovskyi}, pdfauthor={Iurii Pigovskyi}
}
\urlstyle{same}

% ---------- helper macros ----------
% \entryrow{Name (bold, may contain \href)}{Dates}
\newcommand{\entryrow}[2]{%
  \noindent\begin{tabularx}{\linewidth}{@{}X@{\hspace{8pt}}r@{}}
  \textbf{#1} & \colorbox{navyfaint}{\footnotesize\textcolor{navy}{\,#2\,}}\\
  \end{tabularx}%
}
% \role{Job title / degree, italic gray}
\newcommand{\role}[1]{{\itshape\footnotesize\color{gray1}#1}\par\vspace{0.5pt}}
% \techline{Technology list, shaded full-width box}
\newcommand{\techline}[1]{%
  {\setlength{\fboxsep}{2.2pt}%
  \vspace{1pt}\noindent
  \colorbox{gray2}{\parbox{\dimexpr\linewidth-2\fboxsep\relax}{\footnotesize\color{gray1}\textbf{Technologien: }#1}}}
  \vspace{1.5pt}\par
}

\begin{document}

% ================= HEADER =================
\begin{tcolorbox}[
  colback=navyfaint, colframe=navyfaint, boxrule=0pt, arc=3pt,
  width=\linewidth, left=16pt, right=16pt, top=7pt, bottom=7pt,
  before skip=0pt, after skip=7pt
]
{\fontsize{22}{24}\selectfont\bfseries\textcolor{navy}{Iurii Pigovskyi}}\\[1pt]
{\itshape\color{gray1}(Yuriy Pigovsky)}\\[5pt]
{\large\bfseries Java-, Scala- und Python-Softwareentwickler, Softwarearchitekt}\\[5pt]
Dortmund, Deutschland \ \textbar\ \ geb. 11.02.1983 (43 Jahre)\\
\href{mailto:pigovsky@gmail.com}{pigovsky@gmail.com} \ \textbar\ \ WhatsApp / Mobil: +49\,170\,2906351\\
Verfügbar für: unbefristete Festanstellung in Dortmund / Hybrid / Remote\\
Aufenthalts-/Arbeitsstatus: Blaue Karte EU (Blue Card)
\end{tcolorbox}

% ================= ZUSAMMENFASSUNG =================
\section*{Zusammenfassung}
Seit 2004 in der Softwareentwicklung tätig -- insgesamt 22 Jahre Erfahrung. Seit 2015 (11 Jahre) spezialisiert auf Entwurf und Entwicklung verteilter Systeme in Java/Scala und Python:
\begin{itemize}
  \item Java + Scala -- 12 Jahre Erfahrung
  \item Python -- 9 Jahre Erfahrung
  \item Neben eigener Entwicklung: Koordination der Teamarbeit und Durchführung von Code-Reviews (PRs).
  \item Erfahrung in wissenschaftlicher Forschung: mathematische Modellierung, numerische Methoden, maschinelles Lernen.
\end{itemize}

% ================= AUSBILDUNG =================
\section*{Ausbildung}
\entryrow{Nationale Universität „Lviv Polytechnic“}{11/2005 -- 10/2008}
\role{Promotion (PhD) in technischen Wissenschaften}
\vspace{2pt}

\entryrow{Westukrainische Nationale Universität (ehem. Ternopil Academy of National Economy)}{09/1999 -- 07/2004}
\role{Master in Wirtschaftskybernetik}

% ================= FÄHIGKEITEN =================
\section*{Fähigkeiten / Kompetenzen}
\begin{tabularx}{\linewidth}{@{}>{\bfseries}p{4.4cm}X@{}}
Programmiersprachen \& Technologien & Scala (Akka HTTP/Pekko, Specs2, Mockito, Play Framework, sbt); Python (FastAPI, Django, Celery); Java, Kotlin (Spring Boot, Android SDK, JUnit, Gradle, Maven); Bash \\[5pt]
Datenbanken & SQL -- MariaDB, H2, PostgreSQL, SQLite; NoSQL -- Cassandra, Redis; Event-Logging \& Message Queues -- Kafka, RabbitMQ \\[5pt]
Metriken & OpenTelemetry, Prometheus \\[5pt]
Betriebssysteme & Linux, Windows \\[5pt]
Versionsverwaltung & Git (fundierte Kenntnisse) \\[5pt]
Deployment & Docker, Kubernetes (k8s), GCP, GKE \\
\end{tabularx}

% ================= BERUFSERFAHRUNG =================
\section*{Berufserfahrung}

\entryrow{Fachhochschule Dortmund}{07/2024 -- heute}
\role{Wissenschaftlicher Mitarbeiter, Senior Softwarearchitekt \& Distributed Systems Engineer}
\begin{itemize}
  \item \href{https://www.fh-dortmund.de/microsite/smartedgelab/projekte/emulate.php}{Metaprojekt „Emulate“}
  \item Beratung von Telekommunikationsbetreibern; Implementierung von Features und Fixes für deren Business-Support-Lösungen
  \item \href{https://facemap-de.xyz/}{facemap-de.xyz} -- soziales Netzwerk
  \item SeQaM -- Plattform für Service-Quality-Management mit minimaler Ende-zu-Ende-Latenz
\end{itemize}
\techline{Python (FastAPI), Scala (Pekko), Java (Spring Boot), OpenTelemetry, Prometheus, ClickHouse DB, Redis, RabbitMQ, KVM, Docker, k8s, Helm, Multiagentenentwicklung}

\entryrow{\href{https://www.qvantel.com}{Qvantel}}{12/2017 -- 06/2024}
\role{Senior Scala-/Python-Entwickler}
\begin{itemize}
  \item BSSAPI -- Business Support Solution für Telekommunikation (Scala, Akka, Pekko, TDD)
  \item DBSS -- Digital Business Support Solution für Telekommunikation (Python, Django); zusätzlich PR-Reviews und Koordination des Teams
  \item Services und Skripte zur Unterstützung des Message Manager
  \item bssapi-v1-to-v0-adapter -- Microservice zur Konvertierung von BSSAPI-Protokollversionen
\end{itemize}
\techline{Scala 2.11/2.12, Akka, Akka HTTP, Akka Stream, Python 3.7, Django, Kafka, RabbitMQ, Cassandra, Redis, MariaDB, Docker, Consul, Marathon, Mesos, TDD}

\entryrow{\href{https://scalhive.com}{ScalHive}}{12/2016 -- 12/2017}
\role{Scala-Entwickler}
\begin{itemize}
  \item CRM -- Automatisierung von Digital-Marketing-Workflows
  \item GeoFence -- Erfassung von Marketingstatistiken
\end{itemize}
\techline{Scala 2.11/2.12, Akka, Akka HTTP, Akka Stream, PostgreSQL}

\entryrow{Ecodery}{06/2015 -- 12/2016}
\role{Java-/Scala-Entwickler}
\begin{itemize}
  \item Chativity (Buzzchat) -- Android-App
  \item El Chatto -- Chatbot
\end{itemize}
\techline{Java, Kotlin, Android SDK, JUnit, TDD, Scala, Akka, Akka HTTP}

\entryrow{Spilna Sprava}{06/2014 -- 06/2015}
\role{Java-Android-Entwickler}
\begin{itemize}
  \item Mon Dentist -- Plattform zur Kommunikation zwischen Zahnärzten und Patienten
\end{itemize}
\techline{Java, Android SDK, JUnit, TDD}

\entryrow{Wissenschaftliche Arbeit}{09/2004 -- 12/2008}
\role{Verfassen und Verteidigung der Dissertation sowie wissenschaftlicher Publikationen}
\begin{itemize}
  \item \href{https://ena.lpnu.ua/items/c882b7c6-b45d-4b1f-87ca-0a9938085e6c}{„Mathematical models of fuzzy processes in Monod-Iyerusalimskii systems and methods of their identification“}
\end{itemize}

\entryrow{Freelance}{06/2004 -- 06/2014}
\role{PHP-Entwickler}
\begin{itemize}
  \item CMS mit erweiterten SEO-Funktionen
\end{itemize}
\techline{PHP, JavaScript, CodeIgniter}

\entryrow{Westukrainische Nationale Universität}{09/2004 -- 12/2017}
\role{Dozent (Senior Lecturer) -- Fächer: Betriebssysteme, Java-Technologie}

% ================= ERFOLGE =================
\section*{Erfolge}
\begin{itemize}
  \item Mitwirkung am Entwurf und der Implementierung von Prepaid- und Postpaid-Zahlungssystemen, die aktuell im großflächigen produktiven Einsatz bei mehreren großen Mobilfunkbetreibern in verschiedenen Ländern stehen.
  \item Integration des Message Manager bei zwei großen Mobilfunkbetreibern inklusive zahlreicher unterstützender Skripte.
  \item Erfahrung sowohl mit DBSS (Python) als auch mit BSSAPI (Scala) war bei der Entwicklung neuer Features maßgeblich hilfreich.
  \item Implementierung eines k8s-Cluster-Monitors und -Recommenders zur Verbesserung der Servicequalität für geschäftskritische Anwendungen.
\end{itemize}

% ================= ZERTIFIKATE =================
\section*{Zertifikate}
\begin{itemize}
  \item \href{https://coursera.org/share/2ebc59f06b3331339859133d6571aafa}{Unsupervised Learning, Recommenders, Reinforcement Learning}
  \item \href{https://coursera.org/share/5bd5fb3081e2dc3503a12073b73c5bee}{Advanced Learning Algorithms}
  \item \href{https://coursera.org/share/1449f4013aaff33a5da5d8b459521d4d}{Supervised Machine Learning: Regression and Classification}
  \item \href{https://coursera.org/share/15d40eb8e6f00ecbb3d0141bada2779e}{Big Data Analysis with Scala and Spark}
\end{itemize}

% ================= SPRACHEN =================
\section*{Sprachen}
\begin{tabularx}{\linewidth}{@{}>{\bfseries}p{4.4cm}X@{}}
Ukrainisch & Muttersprache \\[4pt]
Englisch & tägliche Arbeitssprache \\[4pt]
Deutsch & \href{https://results.telc.net/vb?credential=telc-LcFTx7e}{B1 (telc-Zertifikat)} \\[4pt]
Niederländisch & Grundkenntnisse im Verstehen, Sprechen und Lesen \\
\end{tabularx}

\end{document}
